\documentclass[11pt]{article}

\usepackage[margin=1in]{geometry}
\usepackage{amsmath, amssymb}
\usepackage{booktabs}
\usepackage{enumitem}
\usepackage{graphicx}
\usepackage{hyperref}
\usepackage{longtable}
\usepackage{tabularx}
\graphicspath{{../../artifacts/showcase/plots/}{../../artifacts/bayesian_walkthroughs/plots/}}

\title{Kinematic Classifier Workflow and Math}
\author{kinematic-classifier-sandbox}
\date{\today}

\begin{document}

\maketitle

\begin{abstract}
This document is a practical walkthrough of the repository's kinematic-classification workflow. Its purpose is narrower than the full methodology paper: it explains what flows through the system, how corpus, feature, evidence, and posterior layers connect, and which equations govern the main study loop. The current 1D witness suite is used as a readable proof of workflow layers, not as the final scientific endpoint. The workflow language combines corpus governance, Bayesian posterior updating, and evidence-provider comparisons using standard statistical references as the mathematical baseline rather than claiming a new end-to-end theory \cite{bishop2006,kalman1960,brier1950}.
\end{abstract}

\section{Purpose}
The repository can be read at two levels:
\begin{enumerate}[leftmargin=1.5em]
  \item as a generic study workflow for kinematic classification, and
  \item as a current 1D witness suite that proves parts of that workflow.
\end{enumerate}

This document focuses on the workflow itself. The main question is:
\begin{quote}
\textit{How does a declared kinematic classification study move from corpus design to a promotion decision under explicit mathematical contracts?}
\end{quote}

\section{End-To-End Flow}
The shortest accurate summary of the repository is
\begin{align}
\text{study declaration}
&\rightarrow
\text{corpus generation or selection}
\rightarrow
\text{class-validity and coverage audits}
\notag \\
&\rightarrow
\text{feature extraction or model-based evidence}
\rightarrow
\text{posterior updates}
\notag \\
&\rightarrow
\text{evaluation, stress tests, and promotion decision}.
\end{align}

The central study object is
\begin{equation}
s = (\mathcal{D}, f, \mathcal{C}, m, \pi, b),
\end{equation}
where
\begin{align}
\mathcal{D} &:= \text{corpus or selected dataset slice}, \\
f &:= \text{feature set}, \\
\mathcal{C} &:= \text{class set or class pair}, \\
m &:= \text{classifier or filter family}, \\
\pi &:= \text{prior regime}, \\
b &:= \text{optional backend or dynamics family}.
\end{align}

The workflow then has two generic engines:
\begin{align}
\text{Corpus Explorer} &: (o, b) \mapsto \mathcal{D}^{\star}, \\
\text{Study Candidate Evaluator} &: (\mathcal{D}^{\star}, f, \mathcal{C}, m, \pi, b) \mapsto d,
\end{align}
with decision
\begin{equation}
d \in \{\text{promote}, \text{revise}, \text{reject}, \text{defer}\}.
\end{equation}

The current strongest supported corpus-evaluation workflow is the default generated/common-study path. The selected-corpus rerun path is also strong because the chosen generated packet is pushed back through adequacy, leakage, and regression checks. An adapter/contract layer exists for provided corpus sources, but it is not yet a full arbitrary-corpus adequacy pipeline, so the repo should not be described as providing full generic corpus evaluation for any provided corpus.

Figure~\ref{fig:study_object_map} is the compact front-door diagram for that object.

\begin{figure}[htbp]
\centering
\setlength{\fboxsep}{8pt}
\fbox{
\begin{minipage}{0.92\textwidth}
\centering
\begin{tabularx}{\textwidth}{@{}XcXcX@{}}
\textbf{Declared Study} &
\(\rightarrow\) &
\textbf{Corpus Engine} &
\(\rightarrow\) &
\textbf{Decision Engine} \\
\(s = (\mathcal{D}, f, \mathcal{C}, m, \pi, b)\) &
&
\((o, b) \mapsto \mathcal{D}^{\star}\) &
&
\((\mathcal{D}^{\star}, f, \mathcal{C}, m, \pi, b) \mapsto d\) \\
\end{tabularx}

\vspace{0.7em}

\begin{tabularx}{\textwidth}{@{}>{\raggedright\arraybackslash}X>{\raggedright\arraybackslash}X>{\raggedright\arraybackslash}X@{}}
\textbf{\(\mathcal{D}\)} corpus slice or selected dataset &
\textbf{\(f\)} feature set and evidence summary family &
\textbf{\(\mathcal{C}\)} class set or class pair \\
\textbf{\(m\)} classifier or filter family &
\textbf{\(\pi\)} prior regime &
\textbf{\(b\)} backend, sensor, or dynamics family \\
\end{tabularx}
\end{minipage}
}
\caption{The repository's front-door study object. The left box declares what is being claimed, the middle engine selects a corpus that can test the claim honestly, and the right engine decides whether the declared study is promotable.}
\label{fig:study_object_map}
\end{figure}

\section{Workflow Stages}
\subsection{Stage 1: Declare the Study}
The workflow starts by declaring what is being tested rather than jumping directly to a model run. That declaration fixes:
\begin{itemize}[leftmargin=1.5em]
  \item the classes to separate,
  \item the kinematic feature family to use,
  \item the evidence-provider family to compare,
  \item the prior regime,
  \item and the corpus requirements needed to make the claim honest.
\end{itemize}

This separation matters because poor results can have multiple causes:
\begin{equation}
\text{failure} \neq \text{classifier failure only}.
\end{equation}
A bad result can indicate inadequate corpus coverage, weak features, ambiguous class definitions, prior fragility, or a missing algorithmic rung.

\subsection{Stage 2: Generate or Select a Corpus}
For the current 1D witness suite, a trajectory is represented as
\begin{equation}
\tau_i = \{(t_{ik}, z_{ik})\}_{k=1}^{n_i},
\end{equation}
with latent kinematic state
\begin{equation}
x_{ik} =
\begin{bmatrix}
p_{ik} \\
v_{ik} \\
a_{ik}
\end{bmatrix},
\qquad
z_{ik} = h(x_{ik}) + \epsilon_{ik}.
\end{equation}

In the generic story, the corpus layer defines a distribution over scenarios rather than one fixed benchmark file:
\begin{equation}
\tau \sim p(\tau \mid \theta, g, \eta),
\end{equation}
where \(\theta\) are scenario parameters, \(g\) is a class or motion template, and \(\eta\) captures noise, corruption, or stress settings.

The corpus objective is not merely ``make more data.'' It is to produce a dataset that exercises the intended boundaries:
\begin{align}
Q_{\text{corpus}}
=\;&
w_{\text{balance}} Q_{\text{balance}}
+ w_{\text{pair}} Q_{\text{pair}}
+ w_{\text{feature}} Q_{\text{feature}}
\notag \\
&+ w_{\text{difficulty}} Q_{\text{difficulty}}
- w_{\text{leak}} P_{\text{leak}}
- w_{\text{invalid}} P_{\text{invalid}}.
\end{align}
The exact score terms vary by analysis surface, but the structure is stable: reward coverage and difficulty diversity, penalize leakage and invalid examples.

\subsection{Stage 3: Audit Coverage Before Classifier Claims}
Before comparing classifiers, the workflow checks whether the corpus itself is methodologically usable. The main audit questions are:
\begin{itemize}[leftmargin=1.5em]
  \item Do declared classes remain valid after generation and corruption?
  \item Are feature sets actually excited by the generated trajectories?
  \item Are hard class pairs represented near their intended boundaries?
  \item Is the corpus balanced enough to support downstream interpretation?
  \item Are there obvious covariate leaks that make the task artificially easy?
\end{itemize}

In compact form, the adequacy gate is
\begin{equation}
\text{corpus adequate}
\Longleftrightarrow
\bigwedge_{j=1}^{J} G_j(\mathcal{D}) = \text{pass},
\end{equation}
where each \(G_j\) is a gate such as class validity, balance, feature-set coverage, or leakage.

The canonical artifact path for this stage is a paired adequacy-and-coverage packet: \texttt{corpus\_adequacy\_audit\_v1} establishes the corpus gates themselves, while \texttt{coverage\_report\_v1} records classifier-support and feature-set readiness against the same audited corpus.

\subsection{Stage 4: Convert Observations Into Evidence}
The repository separates raw observations from derived features and from posterior-ready evidence:
\begin{align}
z_k &= \text{raw observation at time } t_k, \\
\phi_k &= \Phi(z_{1:k}, t_{1:k}), \\
y_k &= \text{evidence-provider input at step } k.
\end{align}

The evidence input \(y_k\) depends on the method family:
\begin{equation}
y_k =
\begin{cases}
z_k, & \text{pointwise methods}, \\
\phi_k, & \text{windowed feature methods}, \\
\nu_{k,c}, & \text{Kalman innovation evidence}, \\
(\phi_k, \nu_{k,c}), & \text{hybrid or future combined methods}.
\end{cases}
\end{equation}

The important design rule is that methods are compared through a shared evidence contract:
\begin{equation}
\mathcal{E}(z_{1:k}, \phi_{1:k}, \text{model state}) \rightarrow \{\ell_k(c)\}_{c \in \mathcal{C}},
\end{equation}
where \(\ell_k(c)\) is the class log-evidence contribution at time \(k\).

\subsection{Stage 5: Update the Posterior}
Once evidence is expressed in a compatible form, the posterior update is generic. Let
\begin{equation}
p_{k}(c) := p(c \mid y_{1:k}).
\end{equation}
Then the recursive Bayes update is
\begin{equation}
p_k(c) \propto p_{k-1}(c) \exp(\ell_k(c)),
\end{equation}
with normalized form
\begin{equation}
p_k(c) =
\frac{p_{k-1}(c) \exp(\ell_k(c))}
{\sum_{c' \in \mathcal{C}} p_{k-1}(c') \exp(\ell_k(c'))}.
\end{equation}

In log space, define
\begin{equation}
\alpha_k(c) = \log p_{k-1}(c) + \ell_k(c),
\end{equation}
then
\begin{equation}
p_k(c) = \frac{\exp(\alpha_k(c))}
{\sum_{c' \in \mathcal{C}} \exp(\alpha_k(c'))}.
\end{equation}

For a two-class problem with classes \(c_1\) and \(c_2\), the log-odds recursion is
\begin{equation}
\log \frac{p_k(c_1)}{p_k(c_2)}
=
\log \frac{p_{k-1}(c_1)}{p_{k-1}(c_2)}
+
\ell_k(c_1) - \ell_k(c_2).
\end{equation}
This equation is central to the repo because it exposes prior sensitivity and flip thresholds directly.

\subsection{Stage 6: Add Structure When Static Evidence Is Not Enough}
If pointwise or windowed evidence is insufficient, the workflow adds structure one rung at a time rather than leaping immediately to the most complex filter.

\paragraph{Transition-aware accumulation.}
For switching-aware inference, class posteriors are propagated through a transition matrix \(T\):
\begin{equation}
\tilde{p}_k(c') = \sum_{c \in \mathcal{C}} T_{cc'} \, p_{k-1}(c),
\end{equation}
followed by the evidence update
\begin{equation}
p_k(c') \propto \tilde{p}_k(c') \exp(\ell_k(c')).
\end{equation}

\paragraph{Model-based innovation evidence.}
For Kalman-bank style methods, each class carries a dynamics model and emits an innovation \(\nu_{k,c}\) with covariance \(S_{k,c}\). The class evidence increment is
\begin{equation}
\ell_k(c)
=
-\frac{1}{2}
\left(
\nu_{k,c}^{\top} S_{k,c}^{-1} \nu_{k,c}
+ \log \det S_{k,c}
+ d \log(2\pi)
\right),
\end{equation}
where \(d\) is the observation dimension.

The ladder logic is therefore:
\begin{align}
\text{pointwise}
\rightarrow
\text{windowed}
\rightarrow
\text{sequential Bayes}
\rightarrow
\text{Kalman bank}
\notag \\
\rightarrow
\text{transition-aware filters}
\rightarrow
\text{advanced filters}.
\end{align}
Each rung must fix a demonstrated failure mode of the previous one.

\subsection{Stage 7: Evaluate the Study}
After posterior histories are produced, the evaluator does more than inspect top-line accuracy. It measures:
\begin{itemize}[leftmargin=1.5em]
  \item confusion and pairwise failure localization,
  \item prior sensitivity,
  \item separability of the feature space,
  \item calibration and posterior sharpness,
  \item robustness under stress, corruption, and boundary cases.
\end{itemize}

The evaluation layer asks two different questions:
\begin{align}
\text{Is the feature/class geometry promising?} \\
\text{Does the chosen evidence-provider family realize that promise?}
\end{align}

That is why the workflow includes both static geometry diagnostics and dynamic posterior diagnostics.

\subsection{Stage 8: Promote, Revise, Reject, or Defer}
The final output is a decision, not just a score:
\begin{equation}
d = \Delta(
\text{corpus audits},
\text{separability},
\text{posterior quality},
\text{stress results}
).
\end{equation}

The point of the decision layer is to keep the repository honest. A study should not be promoted because one artifact looked good in isolation. It should be promoted only when the corpus, feature, evidence, and posterior stories are mutually consistent.

\section{Worked Numeric Examples}
This section uses current repository artifacts so the equations above can be checked against real runs rather than only against symbolic notation.

\subsection{Sequential Bayes Example}
The simplest hand-checkable update in the current artifacts is the binary demo recorded in the Bayes step table artifact. It starts with a uniform prior
\begin{equation}
p_0(A) = p_0(B) = 0.5
\end{equation}
and one-step log evidence values
\begin{align}
\ell_1(A) &= -0.2231435513, \\
\ell_1(B) &= -1.6094379124.
\end{align}

The log Bayes factor is therefore
\begin{equation}
\ell_1(A) - \ell_1(B) = 1.3862943611.
\end{equation}
With equal priors, the posterior odds equal the Bayes factor, so
\begin{equation}
\log \frac{p_1(A)}{p_1(B)} = 1.3862943611
\end{equation}
and the normalized posterior is
\begin{equation}
\bigl(p_1(A), p_1(B)\bigr) = (0.8, 0.2).
\end{equation}

The same artifact also contains a real common-study trajectory where the final log posterior odds for
\texttt{constant\_velocity} versus \texttt{constant\_acceleration} reach
\begin{equation}
\log \frac{p(\texttt{constant\_velocity} \mid y_{1:4})}
{p(\texttt{constant\_acceleration} \mid y_{1:4})}
= 0.7005868961,
\end{equation}
giving final posteriors
\begin{equation}
\bigl(0.6683178817,\; 0.3316821183\bigr).
\end{equation}

\subsection{Kalman Innovation Example}
For the Kalman-bank rung, the scalar innovation likelihood for class \(c\) is
\begin{equation}
\ell_k(c)
=
-\frac{1}{2}
\left(
\frac{\nu_{k,c}^2}{S_{k,c}}
+ \log S_{k,c}
+ \log(2\pi)
\right).
\end{equation}

Using the row for \texttt{constant\_velocity\_regular\_0} at step \(4\) from
\texttt{innovation\_history.csv}, the
\texttt{constant\_velocity} model has
\begin{align}
\nu_{k,\mathrm{cv}} &= -0.1727524922, \\
S_{k,\mathrm{cv}} &= 0.1350588971.
\end{align}
Substitution gives
\begin{equation}
\ell_k(\mathrm{cv})
=
-\frac{1}{2}
\left(
\frac{(-0.1727524922)^2}{0.1350588971}
+ \log(0.1350588971)
+ \log(2\pi)
\right)
= -0.028399.
\end{equation}

On the same step, the \texttt{constant\_acceleration} model has
\begin{align}
\nu_{k,\mathrm{ca}} &= 0.4632798272, \\
S_{k,\mathrm{ca}} &= 0.4137106782,
\end{align}
which yields
\begin{equation}
\ell_k(\mathrm{ca}) = -0.737038.
\end{equation}
Those exact values match the corresponding posterior-history row, where the
\texttt{constant\_velocity} model dominates that step.

\subsection{Transition-Matrix Example}
The transition-aware numeric walkthrough
shows the propagation step explicitly. At the first switched step for the target
mode \texttt{braking},
\begin{equation}
\bar{p}_t(\texttt{braking})
= \sum_s p_{t-1}(s)\,T_{s,\texttt{braking}}
= 0.312.
\end{equation}

The contributing terms are
\begin{align}
0.340 \cdot 0.800 &= 0.272, \\
0.660 \cdot 0.060 &= 0.040,
\end{align}
with the remaining source-mode contributions rounding to zero, so the propagated prior mass becomes
\begin{equation}
0.272 + 0.040 = 0.312.
\end{equation}

That walkthrough also records the empirical effect:
\begin{itemize}[leftmargin=1.5em]
  \item static post-switch accuracy: \(0.000\),
  \item transition-matrix post-switch accuracy: \(0.333\).
\end{itemize}
So the key mechanism is not a new emission model, but moving prior mass into a plausible switched mode before the emission term is applied.

\subsection{Corpus Adequacy Example}
The corpus-audit scorecard in
\texttt{corpus\_adequacy\_report.md}
currently reports
\begin{align}
Q_{\text{corpus}} &= 0.643, \\
L &= 0.613, \\
T &= 0.500, \\
G &= 0.023.
\end{align}

The same scorecard shows the positive coverage terms
\begin{align}
B_{\text{class}} &= 1.000, \\
B_{\text{tier}} &= 1.000, \\
B_{\text{covariates}} &= 0.704, \\
E_{\text{feature}} &= 0.557, \\
C_{\text{pair}} &= 0.730, \\
V &= 1.000.
\end{align}

The gate still fails even though feature coverage passes, because one hard pair is too easy:
\begin{equation}
\texttt{constant\_acceleration} \text{ vs. } \texttt{maneuver}
\end{equation}
has observed
\begin{align}
\text{pairwise AUC} &= 1.000, \\
\text{overlap} &= 0.000,
\end{align}
despite being declared \texttt{hard}. This is exactly the sort of evidence the audit is meant to surface: a ``good'' classifier result on that pair would be untrustworthy because the corpus boundary is not yet credible.

\subsection{Advanced-Filter Decision Example}
The escalation gate for stronger switching or nonlinear filters is also evidence-based. The walkthrough in
\texttt{advanced\_filter\_decision\_numeric\_walkthrough.md}
reports
\begin{align}
\Delta_{\mathrm{post-switch}}^{\mathrm{TM-static}} &= 0.102, \\
\Delta_{\mathrm{overall}}^{\mathrm{TM-static}} &= 0.046, \\
\Delta_{\mathrm{post-switch}}^{\mathrm{TM-kalman}} &= 0.151, \\
\Delta_{\mathrm{overall}}^{\mathrm{TM-kalman}} &= 0.123.
\end{align}

The current implemented decisions are therefore
\begin{align}
\text{IMM justified now} &= \text{False}, \\
\text{particle filter justified now} &= \text{False}.
\end{align}
The important methodological point is that the repo does not promote a stronger evidence-provider family merely because it exists. It promotes that family only after a simpler rung fails for a reason the next rung can specifically explain.

\section{Visual Anchors}
The numeric examples above are easier to internalize when paired with the existing artifact plots. Figure~\ref{fig:bayes_transition_visuals} shows the two most direct posterior-oriented examples, and Figure~\ref{fig:governance_visuals} shows the governance layer that decides whether stronger claims are justified.

\begin{figure}[htbp]
\centering
\includegraphics[width=0.48\textwidth]{posterior_timeline.png}
\hfill
\includegraphics[width=0.48\textwidth]{transition_matrix_diagnostics.png}
\caption{Left: posterior evolution for the Bayes walkthrough family. Right: the transition-matrix witness showing how switching-aware propagation changes behavior after mode changes.}
\label{fig:bayes_transition_visuals}
\end{figure}

\begin{figure}[htbp]
\centering
\includegraphics[width=0.48\textwidth]{corpus_adequacy_scorecard.png}
\hfill
\includegraphics[width=0.48\textwidth]{advanced_filter_decision_matrix.png}
\caption{Left: corpus governance is evaluated before classifier claims are promoted. Right: escalation to more complex filters is decision-gated by measured witness evidence rather than by ambition alone.}
\label{fig:governance_visuals}
\end{figure}

\section{Implementation Map}
Table~\ref{tab:workflow_map} ties the abstract stages to current code surfaces and typical artifacts.

\small
\begin{longtable}{@{}p{0.16\textwidth} p{0.25\textwidth} p{0.27\textwidth} p{0.24\textwidth}@{}}
\caption{Current workflow stages, code surfaces, and representative artifacts.}
\label{tab:workflow_map} \\
\toprule
Stage & Main question & Code surfaces & Representative artifacts \\
\midrule
\endfirsthead
\toprule
Stage & Main question & Code surfaces & Representative artifacts \\
\midrule
\endhead
Study declaration & What classes, features, priors, and method family are being claimed? & \texttt{common\_experiment}; study-candidate helpers; manifests in \texttt{experiments/} & study manifests; feature sets; class-pair manifests \\
Corpus generation & What trajectories should be generated or selected? & \texttt{trajectory\_generator}; \texttt{corpus/exploration} & selected corpus manifests; candidate score tables \\
Validity and adequacy & Is the corpus methodologically usable? & \texttt{class\_validity}; \texttt{coverage\_report} & \texttt{corpus\_adequacy\_audit\_v1/}; balance and leakage reports \\
Feature geometry & Are the classes separated in the active feature space? & \texttt{feature\_analysis}; \texttt{pca\_analysis} & \texttt{feature\_analysis\_v1/}; \texttt{pca\_analysis\_v1/} \\
Evidence provider & How is class evidence constructed at each step? & sequential Bayes; Kalman bank; \texttt{transition\_matrix} & posterior histories; innovation tables; transition diagnostics \\
Promotion logic & Is the declared study promotable? & \texttt{validation\_ladder}; reporting surfaces & validation-ladder reports; promotion scorecards \\
\bottomrule
\end{longtable}
\normalsize

\section{What This Document Proves And What It Does Not}
This document proves the current workflow shape and the core mathematical contracts:
\begin{itemize}[leftmargin=1.5em]
  \item corpus, feature, evidence, posterior, and decision layers are explicitly separated,
  \item method families are compared through shared log-evidence and posterior contracts,
  \item the 1D witness suite is a readable proof of workflow layers,
  \item the repository already has places to audit adequacy before classifier claims.
\end{itemize}

This document does not prove:
\begin{itemize}[leftmargin=1.5em]
  \item that the current 1D corpus is final,
  \item that one method family is universally best,
  \item that 3D lift is complete,
  \item or that every conceptual contract is fully generalized across all backends today.
\end{itemize}

The honest current claim is:
\begin{quote}
\textit{The repository already supports a coherent kinematic-classification workflow with explicit mathematical interfaces, and the 1D witness suite is there to prove those interfaces before broader dimensional and backend lift.}
\end{quote}

\section{Artifact Guide}
Table~\ref{tab:artifact_guide} lists the artifact names referenced in this document. The full paths live under the corresponding artifact directories in the repository.

\small
\begin{table}[htbp]
\centering
\caption{Artifact names referenced directly in this workflow note.}
\label{tab:artifact_guide}
\begin{tabularx}{\textwidth}{@{}lX@{}}
\toprule
Artifact name & Purpose in this document \\
\midrule
\texttt{bayesian\_step\_tables.csv} & Hand-checkable and common-study Bayes posterior updates \\
\texttt{innovation\_history.csv} & Per-step Kalman innovations and innovation variances \\
\texttt{posterior\_history.csv} & Per-step posterior values used to confirm the likelihood calculations \\
\texttt{transition\_matrix\_numeric\_walkthrough.md} & Worked switching-aware propagation example \\
\texttt{corpus\_adequacy\_report.md} & Corpus scorecard and failure recommendations \\
\texttt{advanced\_filter\_decision\_numeric\_walkthrough.md} & Evidence-based gate for IMM and particle-filter escalation \\
\bottomrule
\end{tabularx}
\end{table}
\normalsize

\section{References}
\begin{thebibliography}{9}
\bibitem{bishop2006}
C. M. Bishop.
\newblock \emph{Pattern Recognition and Machine Learning}.
\newblock Springer, 2006.

\bibitem{kalman1960}
R. E. Kalman.
\newblock A new approach to linear filtering and prediction problems.
\newblock \emph{Transactions of the ASME}, 82D:35--45, 1960.

\bibitem{brier1950}
G. W. Brier.
\newblock Verification of forecasts expressed in terms of probability.
\newblock \emph{Monthly Weather Review}, 78(1):1--3, 1950.
\end{thebibliography}

\end{document}
